% Part: applied-modal-logics
% Chapter: epistemic-logic
% Section: language-epistemic-logic

\documentclass[../../../include/open-logic-section]{subfiles}

\begin{document}

\olfileid[hi]{aml}{el}{pal}

\olsection{सार्वजनिक घोषणा तर्कशास्त्र}

गतिक ज्ञानात्मक तर्कशास्त्र हमें उन विधियों को निरूपित करने देते हैं जिनसे
समय के साथ, अथवा नई सूचना मिलने पर, कर्ताओं का ज्ञान बदलता है। इनमें से कई
तर्कशास्त्र ज्ञान के परिवर्तनों को सूचनात्मक \emph{घटनाओं} या
\emph{अद्यतनों} के द्वारा निरूपित करते हैं। अद्यतन का सबसे मूल प्रकार
सार्वजनिक घोषणा है, जिसमें कोई सूत्र सत्यतापूर्वक घोषित किया जाता है और सभी
कर्ता उस घटना को एक साथ देखते हैं। इसके लिए हम भाषा को निम्न प्रकार विस्तृत
करते हैं।

\begin{defn}
मान लें कि $G$ कर्ता-संकेतों का समुच्चय है। सार्वजनिक घोषणाओं वाले बहु-कर्ता
ज्ञानात्मक तर्कशास्त्र की मूल भाषा में निम्नलिखित होते हैं:
\begin{enumerate}
  \tagitem{prvFalse}{!!{falsity} का प्रतिज्ञप्तिक अचर~$\lfalse$।}{}
  \tagitem{prvTrue}{!!{truth} का प्रतिज्ञप्तिक अचर~$\ltrue$।}{}
  \item !!{propositional variable}ों का !!{denumerable} समुच्चय: $\Obj
    p_0$, $\Obj p_1$, $\Obj p_2$, \dots
  \item प्रतिज्ञप्तिक संयोजक: \startycommalist
  \iftag{prvNot}{\ycomma $\lnot$ (निषेध)}{}%
  \iftag{prvAnd}{\ycomma $\land$ (संयोजन)}{}%
  \iftag{prvOr}{\ycomma $\lor$ (वियोजन)}{}%
  \iftag{prvIf}{\ycomma $\lif$ (!!{conditional})}{}%
  \item प्रत्येक $a \in G$ के लिए ज्ञान-संकारक $\Knows_a$।
  \item प्रत्येक !!{formula}~$!B$ के लिए सार्वजनिक-घोषणा संकारक $[!B]$।
\end{enumerate}
\end{defn}

सार्वजनिक-घोषणा संकारक बॉक्स संकारक की तरह कार्य करता है, अतः भाषा की हमारी
आगमनात्मक परिभाषा इसी के अनुरूप दी जाती है:

\begin{defn}
  ज्ञानात्मक भाषा की \emph{!!{formula}ें} आगमनात्मक रूप से इस प्रकार
  परिभाषित होती हैं:
  \begin{enumerate}
  \tagitem{prvFalse}{$\lfalse$ एक परमाण्विक !!{formula} है।}{}

  \tagitem{prvTrue}{$\ltrue$ एक परमाण्विक !!{formula} है।}{}

  \item प्रत्येक प्रतिज्ञप्तिक चर $\Obj p_i$ एक (परमाण्विक) !!{formula} है।

  \tagitem{prvNot}{यदि $!A$ एक !!{formula} है, तो $\lnot !A$ एक
  !!{formula} है।}{}

  \tagitem{prvAnd}{यदि $!A$ और $!B$ !!{formula} हैं, तो $(!A \land
  !B)$ एक !!{formula} है।}{}

  \tagitem{prvOr}{यदि $!A$ और $!B$ !!{formula} हैं, तो $(!A \lor !B)$
  एक !!{formula} है।}{}

  \tagitem{prvIf}{यदि $!A$ और $!B$ !!{formula} हैं, तो $(!A \lif !B)$
  एक !!{formula} है।}{}

  \tagitem{prvIff}{यदि $!A$ और $!B$ !!{formula} हैं, तो $(!A \liff !B)$
  एक !!{formula} है।}{}

  \item यदि $!A$ एक !!{formula} और $a \in G$ है, तो $\Knows_a !A$ एक
  !!{formula} है।
  
  \item यदि $!A$ और $!B$ !!{formula} हैं, तो $[!A] !B$ एक !!{formula} है।

  \tagitem{limitClause}{इसके अतिरिक्त कुछ भी !!{formula} नहीं है।}{}
\end{enumerate}
\end{defn}

!!{formula} $[!A] !B$ का अभिप्रेत पठन है, ``$!A$ को सत्यतापूर्वक घोषित किए
जाने के बाद $!B$ सत्य है।'' सार्वजनिक घोषणाओं के संदर्भ में कभी-कभी
सार्वज्ञान की बात करना भी उपयोगी होगा, अतः भाषा में सार्वज्ञान संकारक
~$\CKnows_G !A$ भी सम्मिलित हो सकता है।

\end{document}
